\chapter[Nations: Kinship or Imagination?]{Nations: Ancient Kinship or Modern Shared Imagination?}
\section{A Small Dark-Red Book}
Pick up an object: a passport.

A little wider than your palm, it has a dark cover stamped with a national emblem and lettering. Its first page usually makes a courteous but forceful request: civil and military authorities should permit the bearer free passage and provide help when needed. Beyond lie pages awaiting stamps and your photograph beside a field labeled nationality.

This little book does something extraordinary: it claims to answer whose person you are and where you belong. It does not answer alone. Behind the page, an invisible community of tens or hundreds of millions collectively guarantees your passage. A border officer does not know you but recognizes the book, having something similar in their own pocket.

Here is a surprise: this book is much younger than your grandfather. Broadly, before World War I, most Europeans needed no documents to cross borders. Buy a train ticket from Paris to Istanbul and nobody asked your nationality. Only after war shattered the world did states hurriedly institutionalize proof of national belonging; international conferences standardized passport formats in the 1920s. Humanity had moved around maps for millennia without carrying proof of belonging to a state; today, without this page, you cannot even leave the airport.

The nationality field is stranger still. Legally, it states which country recognizes and protects you and can tax you. Yet many read ancestry, language, and feeling into it: whose descendants you are, what you speak, whom you cheer or mourn. A naturalized athlete prompts furious comments: with unrelated ancestry and faltering Chinese, how can he represent us? A legal document instantly becomes a war over blood and belonging.

We say you are Chinese or he is French daily. What does that mean? Something ancient flowing through blood from ancestors thousands of years ago, an uninterrupted heirloom? Or a comparatively young invention, roughly your great-grandfather's age, jointly produced by schools, newspapers, maps, and armies and collectively believed?

The difference is no word game. It decides who qualifies as us and who remains them, who owns land, and what one person may demand of another in the nation's name, including death. Let us unpack nation.

\section{An Alsatian Classroom, 1871}
Enter a scene you may know from Chinese-language textbooks. Today, look again beyond what the textbook invited you to see.

Prussia and France went to war in 1870, and France suffered disastrous defeat. The 1871 settlement ceded borderlands Alsace and Lorraine to Prussia. Orders soon reached every school: tomorrow, German becomes the sole instructional language and all French lessons cease.

In 1873, Alphonse Daudet turned this into The Last Lesson. Young Franz is late for school, tempted to skip because he has not learned yesterday's participles. Passing the town hall, he sees people around the noticeboard, source of two years' bad news---war, conscription, occupation---and feels uneasy. Entering class reluctantly, he finds unusual silence. Village elders sit at the back, even the former mayor, all solemn; Monsieur Hamel wears his holiday formal coat.

Hamel tells them this is his final French lesson. Tomorrow the new German teacher arrives.

Suddenly French seems easy to Franz, every word understandable. He regrets time watching pigeons and cutting wood instead of studying. Hamel says French is the world's most beautiful language and must never be forgotten. A conquered, enslaved people preserving its language holds a key to the prison door.

At the closing bell, Hamel tries to speak but chokes. Turning to the board, he writes with all his strength, Long live France! Then he leans his head against the wall and gestures: class is over; go.

Before being moved, do some honest work. This is fiction, quietly omitting an inconvenient fact: most ordinary Alsatians in 1871 spoke a German dialect. Alsatian is closely related to German; standard school French was as distant for many children as Franz's unlearned lessons. If language determines nation, Alsatians apparently should have been German.

But they did not want to be. In his famous 1882 Paris lecture, Ernest Renan deliberately used Alsace: people largely of Germanic descent speaking German dialects repeatedly chose France. Where ancestry gave no answer, choice did.

A child weeping over a language he seldom speaks: that classroom holds this chapter's whole question.

\section{Put Both Answers on the Table}
State the conflict without rushing to declare a winner.

\textbf{Answer one: nations are ancient communities of blood.}

Chinese, French, and Yamato name living extended families stretching from antiquity. Shared ancestors, language, and land root them in time. Schools, newspapers, and passports merely depict and register an already ancient tree. Nations grow; they are not manufactured.

\textbf{Answer two: nations are modern shared imaginings.}

The concept is only two or three centuries old. Before the French Revolution, almost nobody understood themselves nationally. A peasant belonged to a village, a lord's tenancy, and Catholicism, without knowing what being French meant. Schools, newspapers, conscription, railways, and maps forged countless villages speaking mutually difficult dialects into one nation. Descendants of the Yan and Yellow Emperors, the Germanic nation, and an unbroken imperial line sound primordial yet have identifiable birthdays, mostly within two centuries.

Feel each answer's weight and danger.

If the first is right, immutable ancestry predetermines national boundaries at conception. Outsiders can never become insiders; loyal immigrants remain foreign; mixed children belong incompletely on both sides. Some of history's worst events followed this road.

If the second is right, is an imagined nation a fraud or falsehood? Millions fight, suffer imprisonment, and die for it. How does falsehood wield such real power? Tell an exile unable to return for decades that their nation was constructed: are the tears constructed too?

The difficulty crosses three lines: historical facts, genuine feelings, and political uses of each answer. Choose intuitively before continuing: does a nation resemble an old tree or a constructed building?

\section{Hear Both Voices Through}
First fully state the older-community case. Modern invention is fashionable in intellectual circles, sometimes so fashionable that opponents are dismissed as relics. That is unfair.

\textbf{Voice one: nations grow rather than being made.}

This side holds at least three piles of evidence.

First, real continuity. A northern Shaanxi farmer and a local ancestor five centuries earlier speak mutually understandable dialects, honor similar ancestral tablets, and eat similar festival food. Jews dispersed among dozens of countries for nearly two millennia without a state maintained shared scriptural language, festivals, and prayer toward Jerusalem. Israel's founding in 1948 continued that unbroken thread. Greeks under almost four centuries of Ottoman rule preserved language and memory through church and school before rising for independence in the 1820s. Such identities long predate modern nation-state machinery; calling them school-made fails chronology.

Second, real emotional depth. You may theoretically call money collective imagination, but telling a displaced refugee their homesickness was constructed changes nothing. Constructed pain hurts; an imagined home remains home. Qu Yuan drowned himself before nation-states existed, yet Chinese readers still feel Li Sao across two millennia. No government's directive can manufacture that resonance.

Third, the opponent's evidence can support this case. Before nationalism, languages, customs, and memories already clustered locally. Later construction did not draw on blank paper but build on existing foundations. Even a powerful construction crew needs ground beneath it.

This is a strong argument. Remember it for later inspection.

\textbf{Voice two: nations were made, and you can still visit the machinery.}

This side makes an itemized tour.

Machine one: school. In 1870s France, Daudet's period, only about half the population spoke French, fewer fluently. Breton, Occitan, Basque, and German dialects made rural travel feel foreign to Parisians. Late-nineteenth-century free compulsory schooling and universal conscription, mixing young men in one-language barracks, forcibly made peasants French over three or four generations. Boundaries changed little, but inhabitants received new software.

Machine two: newspapers. We return to the theory later. For now picture millions of strangers reading the same morning news and cheering or raging together. Us is printed each morning.

Machine three: shared memory. Textbooks, anniversaries, monuments, and military parades tell newcomers what we experienced together. How much togetherness was assembled afterward will come later.

Consider a sharply defined scene. In October 1928, young people met in Batavia, Dutch East Indies, now Jakarta. From hundreds of islands and languages, Javanese, Sundanese, and Minangkabau participants had different ancestry and historical grievances. They ended by pledging one homeland, Indonesia; one nation, the Indonesian nation; and one language, Indonesian.

Notice the strangeness. Indonesia was then merely a geographical name coined by Europeans a few decades earlier. The archipelago had never existed as a whole in any sense, nor under one ancient dynasty. These three ones were invented there, not excavated. Yet they became the banner under which hundreds of millions escaped Dutch colonialism and established a state. False? People really died for it in 1945. Real? Before 1928 it existed in no sense.

Here is nationalism's troubling duality: the same weapon liberates and excludes. Indonesian nationalism organized the colonized against colonizers; Gandhi's imagined us Indians challenged the British Empire; the Chinese nation mobilized four hundred million against Japan. Reverse the machinery and national purity in twentieth-century Europe expels millions and sends them to death camps. In the 1990s Balkans, former neighbors fight across newly drawn ethnic boundaries. The pen writing us can, with different ink, write eliminate them.

Both voices hold compatible facts: continuity is real and modern forging is real. How do their evidence piles form one picture?

\section{Thinking Bridge: Separate Four Words}
Many national arguments discover that words started fighting first. State, nation, ethnic group, and citizenship are routinely treated as one despite naming four things. Separating them is a portable everyday tool.

\textbf{State}: machinery of territorially bounded government, law, army, taxes, prisons. Cold, hard, officially stamped. School registration and identity cards belong here. Ask whether something has legal force.

\textbf{Nation}: people imagining themselves as us, usually claiming a shared political destiny, either a state or a place within one. Identity, not an official seal, sustains it.

\textbf{Ethnic group}: people identifying through common ancestry, language, and customs without necessarily seeking independence. Many minority groups celebrate their festivals and speak their languages without demanding a state.

\textbf{Citizenship}: the legal contract between you and a state, printed in your passport. It may ignore ancestry: an adopted foreign infant becomes a citizen on registration. It may ignore feeling: someone can hold a passport without ever visiting its country.

Once separated, many tangles unravel. Consider examples:

Kurds: tens of millions with strong national identity and their own languages, dispersed among four countries without a state of their own. Nation and state separate.

Switzerland: one state with four official languages---German, French, Italian, Romansh---whose speakers share stable Swiss identity without abandoning linguistic lives. A unified state and plural ethnicities seem no contradiction.

Overseas Chinese: American or Singaporean passports indicate citizenship, while Lunar New Year and Chinese language indicate ethnic identity. Which country's person depends on context; citizenship and ethnicity occupy different ledgers.

Japan: not until 2008 did its parliament formally recognize the Ainu as indigenous. Okinawa, formerly Ryukyu, had an independent kingdom, language, and royal house until nineteenth-century annexation. The supposedly timeless homogeneous nation is largely a story widely told after Meiji. This counterexample is easy to miss: the most natural, since-antiquity national pictures are often recently narrated.

When someone says outsiders must have different loyalties, you can see ethnicity used as proof of allegiance and ancestry as citizenship's gate. When told patriotism demands something, ask whether country means state machinery, national us, or land and its inhabitants. Three meanings yield three different demands. Many arguments stop there when participants discover different subjects.

\section{Two Testimonies Face to Face}
Read two nearly contemporary primary texts, French and Chinese, giving opposed answers.

At the Sorbonne on March 11, 1882, Renan delivered What Is a Nation? Its best-known sentence reads:

\begin{quote}
The existence of a nation is a daily plebiscite.
\end{quote}
He rejects race, language, religion, and geography one by one, replacing them with two elements: shared past memories, especially great collective deeds, and a present will to continue living and achieving together. More coldly, he calls forgetting, even historical error, essential to nation formation: much genuinely unpleasant history must first be forgotten.

Second, Sun Yat-sen's 1924 Guangzhou lectures on the Three Principles of the People. His first nationalism lecture defines nation as the long-term product of five natural forces: blood, livelihood, language, religion, and customs. States arise instead through domination and force. Broadly: nations are natural, states man-made.

Placed together, these collide. Renan makes will and daily renewed choice essential, denying determinative ancestry and language. Sun makes nature essential, ranking blood first: nations grow rather than being chosen.

Set theoretical right and wrong aside to examine their political work. Eleven years after France's defeat, German-speaking Alsace remained German-held. Ancestry and language would make it Germany's forever; daily consent lets France argue Alsatians want to return. Definitions are weapons, not neutral. Sun likewise sought to overthrow Qing and build a new China from Han, Manchu, Mongol, Hui, and Tibetan peoples. He needed an encompassing Chinese nation, emphasizing natural roots to deny arbitrary assembly and affirm inherent unity. Change definitions and territory and allegiance follow.

This does not accuse either thinker of lying. It asks, of every authoritative national definition, what case its speaker was pursuing when speaking. Solemn intellectual concepts often carry contemporary task lists in their pockets.

\section{Where Nations Come From: Three Competing Accounts}
We have enough material to compare explanations. Distinguish verifiable events, such as Renan speaking or Indonesian youth pledging; disputed causes of nationalism's proliferation; participants' lived world, such as Franz feeling French was life itself; and our values about nationalism. We compare causes here.

\textbf{First: ancient communities awaken, the primordialist account.}

Nations root in ancient ancestry, language, religion, and land; modern nationalism awakens an old tree rather than planting one. Jewish millennia, Greek centuries, and stable dialect regions supply continuity evidence. But why did nations awaken successively worldwide in the nineteenth century, with such similar newspapers, rewritten textbooks, monuments, and conscription? Coincidence cannot explain synchrony.

\textbf{Second: industrial society manufactures nations, the modernist account.}

Ernest Gellner offered a sharp explanation. Agriculture needs no nation: a lifelong local peasant can live by dialect without literacy, wishing only that the state leave his grain alone. Industry requires standardized manuals and mobile workers changing cities and occupations, all literate in one written language and common culture. Only states can provide mass standardized education. Thus nations did not create states; states created nations for industrial needs, with schools the central machine. This explains nineteenth-century timing through industrialization. Its limit is coldness: a production-line by-product cannot explain dying for it. Nobody weeps over an instruction manual.

\textbf{Third: new buildings on old foundations, ethnosymbolism.}

Anthony D. Smith proposes a middle route. Modern nations are constructed, as modernists rightly say, but rarely from nothing. They rebuild upon ethnic memories, myths, symbols, and languages. Independent Greece reclaimed selected ancient glory, repeating Athenian democracy while saying less about Spartan slavery. Chinese descent from Yan and the Yellow Emperor similarly edits a figure once among many legendary rulers into a family encompassing over a billion, widely narrated only from late Qing through Republican times. Ancients did not simply lie; each generation re-edits ancestors for present needs. Continuity supplies foundations, modern forging new floors. The middle way's weakness is renewed argument in every case over each contribution's share.

These accounts are three pairs of glasses, not teams demanding fandom. Looking at your textbook beginning in antiquity, they see an old tree's photograph, a production-line blueprint, or a renovated foundation's archive.

\section{Deeper Challenge: Imagined Communities}
Our heaviest theoretical tool arrives with Benedict Anderson's slim 1983 book Imagined Communities, whose title became a field-wide term. He was a historian of Southeast Asia.

His one-sentence definition deserves four-part unpacking: a nation is an \textbf{imagined}, \textbf{limited}, \textbf{sovereign} \textbf{community}.

Imagined does not mean false. You cannot know every member even of the smallest nation, yet carry their image. You may never visit Mohe or meet Pamir herders, yet include them in we Chinese. Holding one another in mind is imagination requiring daily maintenance.

Limited means even the largest nation has boundaries beyond which lie they. None yet embraces humanity as its entire us.

Sovereign means the idea arose amid retreating royal and church authority, when the people's self-rule became legitimacy's new source. Nations carry political demands from birth.

Community is Anderson's deepest insight. However unequal wealth and status become, nations are imagined as equal and comradely. This helps explain why strangers die for strangers. Mill workers and office owners may have opposed interests, yet the anthem imagines them comrades. No theory can bypass this mystery; Anderson at least names it accurately.

How does this enter millions of minds? \textbf{Print capitalism} is his key mechanism. After printing spread through fifteenth-century Europe, publishing needed markets. Too few read Latin, so publishers sold vernacular books, compressing thousands of spoken dialects into a few print languages. Readers entered invisible circles. Newspapers were especially important: hundreds of thousands of strangers read the same news simultaneously and become angry together, each knowing others are doing likewise. Us is reprinted every morning. Novels, railway timetables, and national maps similarly stitch strangers into a shared meanwhile.

Maps and censuses build too. Maps present boundaries as ancient nature, though villagers on either side may remember meaningless lines. Censuses put people in ethnic boxes once fluid and ambiguous; decades of forms turn boxes into walls.

Shared-memory factories work continuously. Nineteenth-century France distributed our ancestors the Gauls in schoolbooks, even to West African and Indochinese colonial children of unrelated ancestry. That classroom is the workshop in action. Turkey's founder Kemal went further in 1928, decreeing Latin letters instead of centuries-old Arabic script. Overnight, Ottoman documents became unreadable to the new generation, cutting the past and manufacturing new Turks at the break. Indonesia chose more cleverly: instead of majority Javanese, its national language adopted Malay, formerly a trading lingua franca, favoring no large group. Standardization is never just language evolving; it is political decision.

Weigh the theory and its limits. It explains modern birthdays behind ancient claims, language and textbooks as battlefields, and strangers dying for strangers. But manufacturing imagination does not determine which imagining is good, nor whether you should love your constructed us. Theory dismantles the machine; your choice before its parts remains yours.

\section{This Question Is Standing Beside You}
This is not merely history. The 1882 question appears around you in dozens of forms.

In sports, a naturalized athlete under a national flag starts arguments. He gives everything for us, so belongs: Renan's chosen affiliation. His blood differs, so he is a mercenary: primordial ancestry as a gate. The Sorbonne's old question wears a new jersey.

In your class chat, an overseas-born newcomer has local parents but hesitant Chinese; a foreign-looking classmate grew up here, speaks better dialect, and outranks you in games. Who belongs more? Different first reactions have one of this chapter's accounts behind them.

Online regional stereotyping---everyone there is stingy or generous---shrinks nation-building to provinces. Draw a circle, declare common character. Our tools expose statistical impressions mistaken for inherited essence and labels for reality, the same error as naturally German Alsatians.

Shared memory is still manufactured in textbook revisions, new monuments, anniversaries, and live parades; now you work in the workshop too. Visa-free passport rankings, dual-nationality arguments, and whether to post on National Day are small change from the daily plebiscite.

Beware simple extremes. Sacred nations beyond questioning deny maintenance to machinery that produced liberation and slaughter alike; you have seen its modern birthday and workshop. Conversely, imagined means worthless ignores money, companies, and laws, all shared imaginings able to kill or save. Construction is neither falsehood nor meaninglessness. Precisely because we imagine together, we can revise together.

\section{Your Question: Who Casts the Vote?}
Return to the small dark-red book.

In Renan's daily plebiscite, daily is more radical than vote. Belonging is not a once-purchased birthright but a choice requiring renewal. Yesterday's ballot is insufficient; today needs another.

Here is a question without a standard answer:

\textbf{What should determine national membership: ancestry, language, residence, or mutual choice between person and community? Who should recognize it: the individual, state, or majority of us?}

Review your evidence first: German-speaking Alsatians choosing France; youths pledging Indonesia into being; West African pupils inheriting Gaulish ancestors; Turks cut from the past by decree; and your two classmates.

Or question the questioner: when asking whether someone is one of us, are they really examining that person's qualifications or the circle's boundaries? People draw the circle. To whom do you want to hand the pen?
